\documentclass[11pt,a4paper]{article}
\usepackage[utf8]{inputenc}
\usepackage[T1]{fontenc}
\usepackage[spanish]{babel}
\usepackage{hyperref}
\usepackage{geometry}
\geometry{margin=2.5cm}
\usepackage{parskip}

\title{080 el libre albedrio no se termina- se ponen cerrojos}
\author{Alejandro Valenzuela Rivera\\\small Universidad Autónoma de Nuevo León}
\date{\today}
\begin{document}
\maketitle

\begin{quote}
\textbf{Sesgo declarado.} Este trabajo se sostiene desde el sesgo declarado de su autor, y no finge neutralidad. Lo que afirma se puede tumbar, y ese es el punto: el falsificador está escrito, no escondido.
\end{quote}

\section*{Resumen}
El libre albedrio no se entrega cerrado ni se puede terminar: es un proceso sin final, y asi se deja. Lo que se pone son cerrojos, para que no nos maten a todos mientras el proceso sigue abierto. Y un cerrojo no es una prohibicion: es una marca con la llave a la vista.

\section{Por qué se sostiene}
Un cerrojo que se cierra sin llave no protege, encierra, y se reconoce porque el que esta adentro no puede abrirlo. Por eso el libre albedrio se deja abierto: lo que se cierra es el acto que no deja a quien preguntar, no el proceso.

\section{De dónde viene}
De la conversación del 2026-10-26 sobre los cerrojos.

\section{El caso}
Ninguna: es la posicion del autor sobre el libre albedrio.

\section{Método}
Un cerrojo que se pone con la llave a la vista, y que el de adentro puede abrir.

\section{Límites, y qué lo tumba}
Un cerrojo que se cierra sin llave, o que solo abre el de afuera.

\section{Lo que queda abierto}
\textquestiondown{}Que cerrojo deja el proceso abierto y cual lo cierra sin llave?

\section*{Cómo se leyó esta tesis}
Esto se escribió en el repositorio man-made-empathy, del que este candado es una pieza: candados/02 el alma y el libre albedrio/080 el libre albedrio no se termina- se ponen cerrojos.md. La máquina que lo compuso lo hizo desde este mismo archivo, y no agregó nada que no estuviera escrito. Si algo de esto es falso, se dice en el mismo sitio, y se corrige con fecha.

\end{document}
